\documentclass[12pt]{article}
\usepackage[utf8]{inputenc}
\usepackage[T1]{fontenc}
\usepackage{hyperref}
\title{Kalmyk-Oirat Language Resources for K-12 Curriculum Development}
\author{Compiled by Claude Sonnet 5.5 (claude-sonnet-5-5) for Andr\'as Kornai}
\date{}
\begin{document}
\maketitle

\section{Introduction}
Kalmyk-Oirat (ISO 639-3 xal, Glottocode kalm1243) is a Mongolic language of the Republic of Kalmykia in the Russian Federation. It is endangered and is taught as a school subject in some schools. A sociolinguistic study from 1999 reported that no more than about 6\% of young people were fluent. This report draws only on what a limited web search returned. No complete K-12 curriculum or named school textbook with publisher details turned up, so any such materials still need to be located.

\section{Primary Educational Resources}
\begin{itemize}
\item An article describing the structure of an electronic Kalmyk-language textbook, presented as the first of its kind in Kalmyk studies: \url{https://nit.tuva.asia/nit/en/article/view/206}
\item An article on the educational portal ``Khalmg Keln'' (Kalmyk language) describing its status and development: \url{https://nit.tuva.asia/nit/en/article/view/110}
\item Mongolia Society Special Papers, Issue 13, \emph{Early Kalmyk Primers and Other Schoolbooks}. It reproduces extracts from 1925--1930 schoolbooks and shows the orthographic reforms Kalmyk schooling has gone through: \url{https://mongoliasociety.org/monsoc-shop-1/p/special-papers-issue-13-early-kalmyk-primers-and-other-schoolbooks}
\item A conversational Kalmyk course taught through English, attributed in search results to applied linguist Liubov Darzhinova. It appears on the 7000 Languages page: \url{https://www.7000.org/kalmyk}
\item Ghilyana Dordzhieva, a Wikitongues fellow, who has run online Kalmyk conversation classes since 2020 and is developing a 21-unit textbook: \url{https://wikitongues.org/fellows/ghilyana-dordzhieva/}
\end{itemize}

\section{Linguistic and Reference Materials}
\begin{itemize}
\item Freelang Kalmyk--English and English--Kalmyk dictionary. It is small, at roughly 1,900 entries: \url{https://www.freelang.net/dictionary/kalmyk.php}
\item Freelang Oirat dictionary page: \url{https://mail.freelang.net/dictionary/oirat.php}
\item Swarthmore Linguistics Kalmyk resource page. It lists online and PDF Kalmyk--Russian dictionaries: \url{https://wikis.swarthmore.edu/ling073/Kalmyk}
\item Indiana University CeLT language information page: \url{https://celt.indiana.edu/portal/Kalmyk}
\item ELDP archive record, ``Comparative studies of Kalmyk and Oirat'': \url{https://eldp.access.preservica.com/dk0516/}
\item The Oxford Guide to the Transeurasian Languages, chapter on the language: \url{https://academic.oup.com/book/41762/chapter/354238642}
\end{itemize}

\section{Religious and Community Resources}
No religious texts or community organisations were confirmed in these searches. The Wikitongues fellowship and the online classes above are community-driven efforts.

\section{Digital and Audio Resources}
The ``Khalmg Keln'' portal and e-textbook, the online conversational course and the Freelang dictionaries are the digital resources found. No audio collections were confirmed.

\section{Curriculum Development Context}
Kalmyk is taught as a subject rather than used as a medium of instruction. Its orthography has been reformed at least six times, and a 2020 \emph{Mongolian Studies} article notes that no textbook of modern Kalmyk for Western learners exists. The listed resources are a starting point only. A curriculum would need textbooks, audio and school-programme documentation from Russian sources, which these searches did not surface.

\end{document}